% Part: sets-functions-relations
% Chapter: sets
% Section: russells-paradox

\documentclass[../../../include/open-logic-section]{subfiles}


\begin{document}

\olfileid[fa-IR]{sfr}{set}{rus}
\olsection{پارادوکس راسل}

اصل گسترش به ما اجازه می‌دهد نماد $\Setabs{x}{\phi(x)}$ را برای
\emph{آن} مجموعه از $x$ها به کار ببریم که~$\phi(x)$ درباره‌شان صادق
است. بااین‌حال، تمام آنچه اصل گسترش \emph{در واقع} مجاز می‌سازد، فقط
اندیشهٔ زیر است. \emph{اگر} مجموعه‌ای وجود داشته باشد که اعضایش دقیقاً
همان اشیایی باشند که ویژگیِ $\phi$ را دارند، \emph{آنگاه} تنها یک چنین
مجموعه‌ای وجود دارد. به بیان دیگر: پس از تثبیت یک~$\phi$، مجموعهٔ
$\Setabs{x}{\phi(x)}$ یکتا است، \emph{اگر وجود داشته باشد}.

اما این گزارهٔ شرطی مهم است!{} نکتهٔ اساسی آن است که هر ویژگی‌ای امکان
\emph{مجموعه‌سازی} را به دست نمی‌دهد. یعنی برخی ویژگی‌ها مجموعه‌ای را
\emph{تعریف نمی‌کنند}. اگر همهٔ ویژگی‌ها مجموعه تعریف می‌کردند، با
ناسازگاری‌های آشکار روبه‌رو می‌شدیم. مشهورترین نمونهٔ این امر پارادوکس
راسل است.

مجموعه‌ها می‌توانند !!{element}s مجموعه‌های دیگر باشند---برای نمونه،
مجموعهٔ توانیِ مجموعهٔ~$A$ از مجموعه‌ها تشکیل شده است. بنابراین
پرسیدن یا بررسی اینکه آیا مجموعه‌ای !!a{element} مجموعه‌ای دیگر است،
معنا دارد. آیا یک مجموعه می‌تواند عضو خودش باشد؟ ظاهراً هیچ‌چیز در
مفهوم مجموعه این امکان را منتفی نمی‌کند. برای نمونه، اگر \emph{همهٔ}
مجموعه‌ها گردایه‌ای از اشیا را تشکیل دهند، ممکن است گمان کنیم که
می‌توان آنها را در یک مجموعه گرد آورد---مجموعهٔ همهٔ مجموعه‌ها. و این
مجموعه، چون خود یک مجموعه است، !!a{element} مجموعهٔ همهٔ مجموعه‌ها
خواهد بود.

پارادوکس راسل هنگامی پدید می‌آید که ویژگیِ !!a{element} خود نبودن، یعنی
\emph{عضو خود نبودن}، را در نظر بگیریم. اگر فرض کنیم مجموعه‌ای از همهٔ
مجموعه‌هایی وجود دارد که !!a{element} خودشان نیستند، چه؟ آیا
\[
R = \Setabs{x}{x \notin x}
\]
وجود دارد؟ معلوم می‌شود که می‌توانیم ثابت کنیم چنین مجموعه‌ای وجود ندارد.

\begin{thm}[پارادوکس راسل]\ollabel{thm:russells-paradox}
	هیچ مجموعه‌ای $R = \Setabs{x}{x \notin x}$ وجود ندارد.
\end{thm}

\begin{proof}
اگر $R = \Setabs{x}{x \notin x}$ وجود داشته باشد، آنگاه
$R \in R$ اگر و تنها اگر $R \notin R$؛ و این یک تناقض است.
\end{proof}

\begin{tagblock}{novice}
\begin{explain}
بیایید این برهان را آهسته‌تر مرور کنیم. اگر $R$ وجود داشته باشد، پرسیدن
اینکه آیا $R \in
R$ است یا نه معنا دارد. فرض کنید واقعاً $R \in R$. اکنون، $R$ به‌صورت
مجموعهٔ همهٔ مجموعه‌هایی تعریف شده بود که !!{element}s خودشان نیستند. پس
اگر $R \in R$، خودِ $R$ ویژگیِ تعریف‌کنندهٔ $R$ را ندارد. اما فقط
مجموعه‌هایی که این ویژگی را دارند در~$R$ هستند؛ ازاین‌رو $R$ نمی‌تواند
!!a{element}~$R$ باشد، یعنی $R \notin R$. ولی $R$ نمی‌تواند نسبت
به~$R$، هم !!a{element} باشد و هم نباشد؛ پس به تناقض رسیده‌ایم.

چون فرضِ $R \in R$ به تناقض می‌انجامد، داریم $R \notin R$. اما این نیز
به تناقض می‌انجامد!{} زیرا اگر $R
\notin R$، آنگاه خودِ $R$ ویژگیِ تعریف‌کنندهٔ $R$ را دارد، و بنابراین
$R$ نیز درست مانند همهٔ مجموعه‌های دیگری که عضو خود نیستند، !!a{element}
$R$ خواهد بود. و باز هم نمی‌تواند نسبت به~$R$، هم !!a{element} نباشد و
هم باشد.
\end{explain}
\end{tagblock}

\begin{digress}
چگونه نظریه‌ای برای مجموعه‌ها بنا کنیم که به پارادوکس راسل گرفتار نشود؛
یعنی ادعای \emph{ناسازگارِ} وجود مجموعهٔ
$R = \Setabs{x}{x \notin x}$ را مطرح نکند؟ باید اصول موضوعی‌ای وضع کنیم
که شرایط بسیار دقیقی برای بیان اینکه مجموعه‌ها چه هنگام وجود دارند (و چه
هنگام وجود ندارند) به دست دهند.

نظریهٔ مجموعه‌هایی که در این فصل طرح شده چنین کاری نمی‌کند. این نظریه
\emph{به‌راستی ساده‌انگارانه} است. فقط می‌گوید مجموعه‌ها از اصل گسترش
پیروی می‌کنند و اگر چند مجموعه
داشته باشیم، می‌توانیم اجتماع، اشتراک و مانند آن را تشکیل دهیم. می‌توان
نظریهٔ مجموعه‌ها را به‌صورتی دقیق‌تر از این بسط داد.
\oliflabeldef{cumul:::part}{آن دقت را برای بخش
\olref[cumul][][]{part} نگه می‌داریم. فعلاً ساده‌انگارانه، اما با احتیاط،
پیش می‌رویم و می‌کوشیم از ناسازگاری‌ها دوری کنیم.}{}
\end{digress}


\end{document}
